\section{Applications and Evidence from Deployed or Experimental Systems}
\label{sec:applications}

Applications are where an agent architecture becomes a claim about the world. A system does not become dependable merely because its model can reason about a domain: it must encode domain evidence in a representation that survives across steps, constrain what that representation may authorize, and verify the resulting artifact or external effect. The same language model can therefore support a low-risk literature assistant, a consequential account workflow agent, or a laboratory controller, while the required state, authority, and evidence differ sharply. Domain labels alone obscure this distinction. We characterize each application by a \emph{representation--action--verification contract}: what persistent object carries evidence, which bounded effect may be produced from it, and which independent mechanism or accountable person determines whether the transition is valid. AppWorld's executable application state checks and $\tau$-bench's joint treatment of policy and final database state illustrate this contract more directly than answer quality alone~\cite{trivedi2024appworld,yao2024taubench}.

Evidence about an application must also be matched to the claim it can support. A controlled benchmark establishes capability only within its declared initial state, tools, permissions, tasks, and oracle; an experimental closed loop shows that a representation can guide a measurable effect within a particular apparatus; an organizational field study can estimate operational outcomes under the deployed intervention; and a product announcement establishes existence, not effectiveness. Peer review can strengthen confidence in any of these designs, but it is not itself an evidence type. The distinction is visible across WorkArena's enterprise software benchmark~\cite{drouin2024workarena}, A-Lab's physical closed-loop system~\cite{szymanski2023alab}, and a field study of an AI assistant used by 5,172 customer support workers~\cite{brynjolfsson2025generativeaiatwork}. The last study reports a 15\% average increase in issues resolved per hour, but people remained responsible for the conversation and could ignore or edit suggestions. It is therefore valuable evidence about human-supervised assistance, not evidence that an autonomous agent can safely make account changes.

The literature remains asymmetrical: controlled agent benchmarks and engineered demonstrations are much more common than field evidence that joins agent configuration, decisions, realized effects, human workload, and downstream outcomes. Rather than smoothing over that gap, the following tables separate the application contract from the evidentiary boundary of each reported system. This preserves breadth while preventing a successful benchmark episode, a laboratory demonstration, and an organizational outcome study from being treated as interchangeable proof.

\begin{table}[!htbp]
\centering
\caption{Application domains as representation--action--verification contracts.}
\label{tab:applications}
\small
\surveytableformat
\begin{tabularx}{\textwidth}{L{2.45cm}Y Y Y}
\toprule
\rowcolor{surveyBlue!28}\textbf{Domain} & \textbf{Evidence-bearing representation} & \textbf{Bounded effect and principal risk} & \textbf{Verifier and authority gate} \\
\midrule
Enterprise knowledge work & Claim--evidence ledger with source, version, access scope, and contradiction state & Cited report or business object update; unsupported claims, stale policy, or access leakage & Entailment and version checks, authorization-aware retrieval, artifact review \\
Customer operations & Case state linked to policy version, customer authorization, and proposed transaction & Order, credit, or service change; incorrect commitment, privacy leakage, or poor handoff & Deterministic eligibility policy, preview, user confirmation, receipt, escalation \\
Software engineering & Issue contract, repository commit, patch, test record, and unresolved obligation list & Reviewable code change; insecure or incomplete patch and hidden regression & Isolated build and tests, static analysis, diff approval, code owner review \\
IT and security operations & Incident state, telemetry references, runbook version, change plan, and rollback point & Bounded remediation; disruption, credential misuse, or action on adversarial telemetry & Read-only diagnosis, simulation or canary, approval, postcondition, rollback and audit \\
Scientific research & Hypothesis, protocol, instrument state, measurement, and analysis provenance & Program, simulation, or physical experiment; invalid inference, untraceable conditions, or physical hazard & Objective score or measurement, protocol checks, safety interlock, domain expert judgment \\
Healthcare and legal support & Case evidence, authoritative source version, recommendation, uncertainty, and decision record & Draft, recommendation, or routed case; high-consequence factual or procedural error & Source validation, abstention rule, qualified professional authority, outcome monitoring \\
Personal productivity & User request, recipients, permissions, preference scope, and pending action record & Message, schedule, file, or application change; wrong recipient or overbroad access & Draft-first preview, confirmation, reversible commit, scoped and editable memory \\
\bottomrule
\end{tabularx}
\end{table}

\begin{table}[!htbp]
\centering
\caption{Representative application evidence and the bounded claims it supports.}
\label{tab:application-evidence}
\small
\surveytableformat
\begin{tabularx}{\textwidth}{L{2.35cm}L{2.65cm}L{3.05cm}L{3.15cm}Y}
\toprule
\rowcolor{surveyBlue!28}\textbf{System or environment} & \textbf{Setting and design} & \textbf{Observable artifact or effect} & \textbf{Claim supported} & \textbf{Claim not supported} \\
\midrule
WorkArena\newline\cite{drouin2024workarena} & Controlled ServiceNow-based knowledge work benchmark & Browser trajectory and task-specific application state & Enterprise software interaction under the benchmark's tasks and permissions & Reliability across live integrations, organizational policies, and user populations \\
TheAgent\-Company\newline\cite{xu2025agentcompany} & Self-contained simulated software company and authored work tasks & Files, code, messages, websites, and task completion & Multi-application capability on consequential workplace-like tasks & Autonomous job performance or organizational outcomes in an operating company \\
SWE-agent~\cite{yang2024sweagent} and SWE-bench~\cite{jimenez2024swebench} & Repository issue resolution system and executable benchmark & Patch, execution trace, and test outcome & Repository navigation and repair against the available tests & Complete maintainability, security, or requirements coverage beyond the oracle \\
OSWorld~\cite{xie2024osworld} and AppWorld~\cite{trivedi2024appworld} & Controlled GUI and API application environments & Trajectory plus terminal application or database state & Stateful cross-application action with executable end state checks & Full live authentication, interface drift, interruptions, and unobserved side effects \\
FunSearch\newline\cite{romeraparedes2024funsearch} & Closed-loop program search with an objective scoring function & Candidate program and measured score & Search improvement where proposals are executable and cheaply ranked & Open-ended scientific validity when no objective evaluator is available \\
Support assistant field study~\cite{brynjolfsson2025generativeaiatwork} & Staggered deployment to 5,172 customer support workers & Issues resolved per hour, quality measures, and worker outcomes & Operational effects of editable AI suggestions in one organizational setting & Autonomous policy decisions, account changes, or effects without human responsibility \\
\bottomrule
\end{tabularx}
\end{table}

\begin{table}[!htbp]
\centering
\caption{Representative application evidence, continued: scientific loops and their observable boundaries.}
\label{tab:application-evidence-science}
\small
\surveytableformat
\begin{tabularx}{\textwidth}{L{2.35cm}L{2.65cm}L{3.05cm}L{3.15cm}Y}
\toprule
\rowcolor{surveyBlue!28}\textbf{System or environment} & \textbf{Setting and design} & \textbf{Observable artifact or effect} & \textbf{Claim supported} & \textbf{Claim not supported} \\
\midrule
A-Lab~\cite{szymanski2023alab} & Seventeen-day closed-loop inorganic materials campaign & Protocol, instrument actions, measurements, and characterized material & Physical loop closure: 353 experiments and 36 of 57 targets realized in the reported apparatus & General laboratory autonomy outside its materials, equipment, and operating constraints \\
ChemCrow~\cite{bran2024chemcrow} & Tool-augmented chemistry tasks and reported demonstrations & Plan, domain tool calls, and task artifact & Value of chemistry-specific tools and structured execution & Safe authority over arbitrary experiments, equipment, or materials \\
Discovery\-World\newline\cite{jansen2024discoveryworld} and ScienceAgent\-Bench\newline\cite{chen2025scienceagentbench} & Simulated discovery and curated data analysis tasks & Environment state, generated program, execution, and research artifact & Hypothesis-driven interaction and analysis where effects are observable & Scientific validity or physical safety beyond the authored tasks and evaluators \\
AI Scientist~\cite{lu2024aiscientist} and SciAgents~\cite{ghafarollahi2024sciagents} & Reported open-ended research workflow demonstrations & Ideas, experiments, manuscripts, or knowledge graph traces & Feasibility of joining multiple research stages in one scaffold & Artifact validity or safe open-ended operation without external expert adjudication \\
\bottomrule
\end{tabularx}
\end{table}

\subsection{Enterprise knowledge and research assistants}

Enterprise research is not simply RAG with a longer prompt. It becomes agentic when the system must choose where to look, preserve partially resolved questions, compare incompatible sources, invoke business applications, and turn the result into a durable work product. The consequential output may be a cited analysis, a completed case file, a policy recommendation, or a proposed update to a system of record. Each requires more than fluent synthesis: the system must show which admissible evidence supports which claim, which principal is allowed to see and use it, which uncertainty remains, and what downstream action the evidence may authorize. The enterprise contract is therefore evidence-to-artifact first and evidence-to-action only through an explicit authority boundary.

The persistent representation is a \emph{research case} linking the request and audience to subquestions, admissible source versions, claim-level support, contradictions, coverage gaps, intermediate artifacts, tool receipts, the evolving deliverable, and its review state. Model notes remain distinct from authoritative sources, and stable identifiers let a reviewer reject one inference without discarding the entire case.

This representation clarifies when agency is warranted. A deterministic query or fixed RAG pipeline is usually preferable when a known authoritative corpus answers a narrow, read-only question and the output schema is fixed. Agentic research earns its additional cost and risk when the source path is not known in advance, evidence spans repositories or applications, contradictions require conditional follow-up, or the requested artifact has acceptance criteria that must be checked. The distinction is functional rather than stylistic: multiple retrieval calls do not establish agency unless their results update persistent case state and govern subsequent choices. When the case cannot reach adequate coverage or authority, escalation or an explicitly incomplete artifact is a correct terminal state.

Workplace benchmarks expose progressively harder parts of this contract. WorkArena evaluates browser interaction with enterprise objects, while WorkArena++ expands to 682 compositional workflows designed to require planning, logical or arithmetic reasoning, retrieval, and contextual understanding~\cite{drouin2024workarena,boisvert2024workarenaplus}. TheAgentCompany broadens the substrate to files, code, messages, internal websites, and coworker interaction; its NeurIPS report finds that the strongest evaluated agent completed 30\% of its authored tasks~\cite{xu2025agentcompany}. These environments do not establish organizational effectiveness, but their progression is informative: as evidence and effects cross applications, a final answer ceases to be an adequate state representation. Identity, task dependencies, intermediate artifacts, and verified application state must survive tool and context boundaries.

Open web research reveals a complementary bottleneck. AssistantBench contains 214 realistic, time-consuming tasks collected from users and domain experts; no evaluated model exceeded 25 points in accuracy, and the paper reports near-zero scores for contemporary web agents~\cite{yoran2024assistantbench}. The result is tied to its 2024 systems and closed-form evaluators, but its failure analysis shows that navigation, planning, and transfer of information across steps cannot be inferred from a model's closed-book knowledge. A research case should therefore retain failed queries, source coverage, extracted fields, and unresolved dependencies, rather than repeatedly reconstructing progress from the conversation transcript.

Evidence collection and report quality must also be separated. ALCE evaluates fluency, correctness, and citation quality and reports that, on ELI5, even its strongest evaluated systems lacked complete citation support half of the time~\cite{gao2023alce}. DeepResearch Bench extends the target to 100 expert-authored research tasks across 22 fields and separately evaluates report quality, effective citation count, and citation accuracy~\cite{du2026deepresearchbench}. The benchmarks differ in task, corpus, evaluator, and system generation, so their scores should not be compared directly. Their shared lesson is structural: displaying plausible citations does not prove coverage, and a well-supported collection of excerpts does not prove that the final synthesis follows from them. Claim-level evidence edges, contradiction state, and coverage criteria are required between retrieval and prose.

Enterprise authorization adds a boundary absent from open web benchmarks. Repository access should be resolved against the requesting principal and intended audience at retrieval time and checked again when an artifact is shared or committed. A source label generated inside the prompt cannot grant trust or permission, and a citation must not become a channel for disclosing restricted content. Action should advance through an explicit ladder: read-only discovery, internal draft, reviewed artifact, approved communication, and finally a bounded system-of-record change. Each promotion binds the exact research case version to its reviewer, audience, policy decision, and receipt; a later source or permission change reopens the affected claims and proposed actions.

Evaluation should compare deterministic lookup, fixed RAG, and agentic research under identical repositories, permissions, tasks, and budgets. Measure source admissibility, claim support, conflict and stale source handling, artifact acceptance, reviewer correction, privacy, latency, and cost. The invariant is \emph{bounded research}: committed claims and downstream actions require admissible, version-bound evidence appropriate for the audience, while contradictions and gaps remain explicit.

\subsection{Customer support and transactional workflows}

Customer support becomes a transactional agent problem when conversation can change customer or business state: an order is cancelled, a credit is issued, a service visit is scheduled, or an account is modified. At that boundary, a fluent resolution is not necessarily a valid transaction. Each effect must follow from what the customer requested, what the system of record verifies, what the current policy permits, and what authority has been delegated. The architectural question is therefore not simply whether an LLM can choose a tool. It is whether conversational evidence can be converted into an authorized, reviewable, and verifiably completed effect without collapsing customer claims, business facts, policy, and permission into one prompt.

The persistent object is a \emph{customer case}, not the transcript. It separates customer statements from authenticated account facts, current policy, eligibility, conflicts, and proposed resolution. A contemplated effect adds a versioned transaction proposal, proposal-bound confirmation, idempotency, receipts, postconditions, privacy scope, and escalation owner. Stable identifiers connect every commitment to the evidence and policy that justified it.

This representation induces a state machine: intake, identity and consent, diagnosis, policy retrieval, eligibility, proposal, preview and confirmation, commit, receipt verification, and closure or escalation. The path may move backward when the customer revises the request or new evidence invalidates an assumption. Confirmation is not a reusable conversational sentiment: it authorizes one displayed proposal against particular policy, price, account, and tool state versions. A material version change expires that authorization and returns the case to proposal or eligibility. Likewise, a successful tool response does not close the case until the system of record exhibits the intended postcondition; timeouts, duplicate callbacks, and partial updates remain explicit transaction states rather than being paraphrased away.

Benchmarks expose different failures along this path. In retail and airline workflows, $\tau$-bench joins multi-turn conversation, policy compliance, tool use, and final database state, and its pass$^k$ analysis shows why one successful episode does not establish repeated reliability~\cite{yao2024taubench}. ToolSandbox adds stateful execution, implicit dependencies, and intermediate milestones; its difficult categories include state dependency and insufficient information, directly motivating explicit prerequisites and unresolved field status in the case~\cite{lu2025toolsandbox}. $\tau^2$-bench then gives both the user and the agent tools in a shared telecom environment. Moving from no-user operation to this dual control setting reduced pass$^1$ by 18\% for GPT-4.1 and 25\% for o4-mini in the reported experiment, showing that coordination over separately observed state is itself a reliability burden~\cite{barres2025tau2bench}.

Business process breadth introduces another pressure. CRMArena-Pro evaluates nineteen expert-validated sales, service, and configure--price--quote tasks spanning business-to-business and business-to-customer settings. Its leading evaluated agents achieved about 58\% success on single-turn tasks and about 35\% on multi-turn tasks; the study also reports near-zero inherent confidentiality awareness and a task performance cost from targeted confidentiality prompting~\cite{huang2025crmarenapro}. These results are not directly rankable against $\tau$-bench, ToolSandbox, or $\tau^2$-bench because the tasks, models, interaction protocols, and evaluators differ. Collectively, however, they locate the same systems problem: once evidence and authority are distributed across people, policies, and applications, the conversation alone is too weak a control state.

Field evidence provides a useful but narrower comparator. In a staggered deployment to 5,172 customer support workers, an AI assistant increased issues resolved per hour by 15\% on average; workers remained responsible for the interaction and could edit or disregard the generated suggestions~\cite{brynjolfsson2025generativeaiatwork}. This supports the operational value of suggestion-level assistance, especially as a complement to human judgment. It does not establish that an autonomous system may safely determine eligibility or commit account changes. Evaluation and deployment should therefore compare an assistance baseline, a proposal-only agent, and a transaction-capable agent on the same cases, so that the incremental value and risk of delegated authority are visible rather than attributed to generation quality.

Authority should rise through an explicit effect ladder: read-only explanation, draft recommendation, reversible low-impact change, customer-confirmed transaction, and regulated, high-value, or exceptional action. The model may assemble evidence and propose a transition, but deterministic policy and transaction services outside the model should enforce identity, eligibility, amount limits, separation of duties, confirmation requirements, idempotency, and rollback rules. Each rung has a narrower permitted action set and a stronger verifier. Uncertainty must lower authority: missing identity evidence, contradictory account state, an expired policy, or a mismatched confirmation yields clarification, no action, or escalation rather than a best-effort commitment.

Escalation is consequently a valid terminal transition, not a failure to contain the customer. A useful handoff is a structured transfer rather than a transcript dump: it states the customer's goal, verified facts, applicable policy, proposed or attempted actions, receipts and errors, unresolved questions, urgency, consent, and the decision required from the next owner. The same case can then survive channel changes and team boundaries without forcing the customer or employee to reconstruct the episode. Handoff completeness, customer effort, repeat contact, and time to accountable resolution are therefore part of agent quality alongside automation rate.

Evaluation should measure resolution, policy compliance, unauthorized or duplicate effects, confirmation validity, privacy, rollback, handoff, workload, latency, cost, and repeated consistency under identity, policy, state, and tool faults. The invariant is \emph{authorized resolution}: every effect derives from authenticated case evidence, current policy, delegated authority, proposal-bound consent when required, and a verified receipt. Otherwise the system withholds, rolls back, or escalates.

\subsection{Software engineering}

Software engineering makes the evidence-to-action problem unusually concrete. The input is rarely a complete specification: an issue, stack trace, feature request, or review comment provides partial evidence about an intended change. The consequential output is not generated code but a new repository state that may be reviewed, merged, deployed, and depended upon. A patch is therefore valid only when it is tied to a pinned codebase, an explicit change contract, observable failure and success conditions, and authority over the branch or release target. Executability makes verification possible, but it does not make correctness automatic; tests observe selected behavior and cannot by themselves supply missing requirements, architectural judgment, or ownership.

The persistent \emph{change case} binds request and acceptance criteria to a pinned repository, smallest failing reproducer, observed code and test evidence, causal hypotheses, affected interfaces, patch versions, checks, unresolved obligations, review decisions, and any deployment or rollback receipt. Agent hypotheses remain distinct from repository facts, and every observation is version-bound.

This object supports a disciplined state machine: clarify, check out a pinned revision, reproduce, localize, form and test a causal hypothesis, propose a patch, run the narrowest decisive check, expand to regression and nonfunctional checks, request review, commit, and only then stage or deploy. A failed transition returns evidence to the case rather than encouraging unrecorded trial and error. The minimal reproducer is especially important: without one, a passing post-change test cannot show that the patch addressed the reported failure. The case should also separate existing tests, agent-authored tests, and independent evaluation tests. Allowing the same model to infer the requirement, generate the patch, and define the only oracle creates circular evidence even when every test is green.

Repository benchmarks isolate successive parts of this contract. SWE-bench pairs real GitHub issues with repository snapshots and executable tests, establishing patch generation against a fixed task and oracle~\cite{jimenez2024swebench}. SWE-agent shows that the agent--computer interface--how the system searches, edits, and observes the environment--materially changes results~\cite{yang2024sweagent}. Agentless supplies an important simpler baseline: its constrained localization, repair, and patch validation pipeline was competitive with substantially more autonomous systems on SWE-bench Lite and also exposed problems in issue descriptions and tests~\cite{xia2024agentless}. The comparison does not show that agency is unnecessary in general. It shows that adaptive tool choice must earn its complexity against a staged, inspectable workflow on the same repository tasks, models, and budgets.

Broader benchmarks change what the case must preserve. SWE-Lancer contains more than 1,400 paid freelance tasks worth approximately \$1 million in aggregate, spanning small bug fixes, feature implementations, and engineering management choices; its individual contributor tasks use end-to-end tests independently checked by experienced engineers~\cite{miserendino2025swelancer}. SWE-EVO moves from isolated issues to 48 evolution tasks across seven mature projects, averaging 21 changed files and 874 tests per task. In its reported comparison, GPT-5 with OpenHands resolved 21\% of these evolution tasks versus 65\% on the single-issue SWE-bench Verified setting~\cite{thai2025sweevo}. The task sets and scores are not directly rankable, but the widening gap is informative: long-horizon work requires the change case to preserve cross-file obligations, intermediate compatibility decisions, and partial progress that a final patch-and-pass label hides.

Executable tests are consequently evidence, not an infallible specification. An OpenAI audit examined 138 SWE-bench Verified tasks that o3 did not solve consistently and judged 59.4\% to contain material test design or problem description issues, including tests that were narrower or wider than the stated task~\cite{openai2026swebenchaudit}. This targeted audit does not invalidate every benchmark result, and its denominator is not the full suite. It demonstrates why a merge decision cannot reduce to a hidden-test bit: reviewers need the issue-to-acceptance mapping, the behavior exercised by each test, compatibility and security checks, and explicit residual risk. Independent validation should consume the proposed diff and evidence bundle; it should not silently rewrite the patch, relax the acceptance criteria, or treat model confidence as proof.

Field evidence again supports assistance more strongly than autonomous change authority. Across randomized trials at Microsoft, Accenture, and an anonymous Fortune 100 company, access to an AI coding assistant increased completed tasks by 26.08\% in the pooled estimate over 4,867 developers, with substantial uncertainty and variation across experiments~\cite{cui2026highskilledwork}. Developers and organizational controls still owned review, integration, and release. A matched study of more than 100,000 GitHub developers reports that activity gains associated with autonomous coding agents attenuated markedly from commits to projects and releases~\cite{demirer2026writingcode}. The designs estimate different quantities and should not be compared as a leaderboard. Together they motivate measuring accepted and shipped outcomes, review burden, and downstream defects rather than equating generated code or commit volume with delivered value.

Authority should rise through a software effect ladder: read-only diagnosis, localized recommendation, draft diff, tested pull request, approved merge, staged deployment, and production release. Repository and infrastructure services--not the language model--should enforce branch protection, code ownership, credential scope, required checks, deployment policy, and rollback. Execution belongs in an isolated environment with restricted network access, ephemeral credentials, resource limits, and explicit controls on generated scripts and dependencies. High-risk migrations, authentication changes, security controls, and irreversible data operations require stronger independent review and rehearsal. When the case lacks a reproducer, acceptance criterion, required secret, owner, or safe test environment, clarification or handoff is the correct outcome.

Evaluation should compare an assistant, a fixed repair pipeline, and an adaptive agent under identical issues, revisions, tools, models, and budgets. Measure reproduction, localization, acceptance coverage, regressions, security, unauthorized changes, review burden, merge or rollback, escaped defects, latency, and cost. The invariant is \emph{evidence-complete change}: merge or deployment requires a pinned request and repository, reproduced behavior, explicit criteria, diff, independent checks, accountable approval, and receipt.

\subsection{Computer use and personal productivity agents}

Computer use becomes agentic when a language model can alter a user's digital environment rather than merely explain what to click. Email, calendars, documents, browsers, and files make this capability immediately useful, but they also expose the distinction that governs the broader survey: fluency at a screen is not evidence that the resulting effect is correct. Intent, identity, application state, recipients, permissions, and consequences are distributed across interfaces and services. A credible productivity agent must therefore transform a request and heterogeneous observations into an explicit state that can support bounded action; a transcript and the pixels currently visible on screen are insufficient substitutes.

The persistent \emph{personal workflow case} binds the user's goal and deadline to principal, account, application, authoritative service state, versioned source objects, recipients, visibility, plan, and unresolved conflicts. Preferences retain source, scope, expiry, and editability: they may guide a draft but cannot authorize a new recipient, disclosure, purchase, or changed fare.

Pending actions carry typed arguments, preconditions, impact, reversibility, idempotency, expiry, and required confirmation. Exact previews bind approval to attempts, receipts, and postconditions. Cursor position is not identity, focus is not authority, webpage text is not a user instruction, and visual success is not proof of backend state.

The resulting workflow is a state machine: intake and clarification; identity and permission inspection; state gathering; planning; preparation of a private draft; preview and confirmation; commit; authoritative verification and reconciliation; and closure, rollback, or escalation. Before commit, the agent re-queries material state because a recipient, event, file version, price, login, or sharing policy may have changed while it was reasoning. Confirmation is version-bound: it expires if the content, recipient, time, application state, or consequence shown to the user changes. After commit, the system verifies a receipt and the relevant service state instead of inferring success from a disappearing button or a plausible confirmation page.

Existing benchmarks illuminate different parts of this problem rather than forming a directly comparable leaderboard. WebArena evaluates end-to-end tasks on functional websites and reported a 14.41\% success rate for its strongest GPT-4-based agent against 78.24\% for humans~\cite{zhou2024webarena}. OSWorld spans 369 tasks across real desktop operating systems and applications; its initial evaluation reported 12.24\% for the best model compared with 72.36\% for humans~\cite{xie2024osworld}. AndroidWorld adds 116 dynamically instantiated tasks across 20 Android applications and reported 30.6\% for its strongest evaluated agent, while showing sensitivity to task variation and weaker transfer across platforms~\cite{rawles2024androidworld}. AppWorld instead exposes 457 APIs across nine simulated everyday applications and evaluates 750 tasks with state-based tests that can recognize multiple valid solutions and detect collateral changes~\cite{trivedi2024appworld}. The substrates, task distributions, evaluators, models, and dates differ, so their percentages should not be ranked. Together they show why both interaction correctness and resulting application state belong in the evidence chain.

Longer workflows make persistent state more consequential. OSWorld 2.0 contains 108 workflows whose median human completion time is approximately 1.6 hours; Claude Opus 4.7 averaged 318 tool calls, roughly an order of magnitude more than agents in the earlier benchmark. Yet the best reported primary binary completion rate was 20.6\%, with failures involving lost constraints, missed information that arrived mid-task, ungrounded guesses, skipped verification, and hidden application state~\cite{yuan2026osworld2}. These are not only perception errors. They are failures to preserve and reconcile the workflow case as execution unfolds.

Capability and safety must also be measured separately. ClawsBench evaluates 44 tasks across stateful simulations of email, chat, calendar, documents, and storage. Across six models, four agent harnesses, and 33 conditions, fully scaffolded configurations achieved 39--64\% task success while still taking unsafe actions in 7--33\% of runs; the leading harnesses did not exhibit a consistent safety ordering~\cite{li2026clawsbench}. An agent can therefore finish the requested task and still expose the wrong file, contact the wrong person, or make an unauthorized persistent change. Conversely, refusing or failing to complete a task is not evidence of principled safety. Effect correctness, policy compliance, and useful completion require separate denominators.

Field evidence currently supports assistance more strongly than autonomous effect authority. In a six-month randomized experiment across 66 firms and 7,137 knowledge workers, 80\% of treated workers who used a generative AI tool spent about two fewer hours on email per week, while individual access did not produce a detectable change in the quantity or composition of their tasks~\cite{dillon2025shiftingwork}. The result demonstrates practical value from assistance, not that an agent should send mail, alter calendars, or share files without review. Evaluation and deployment should therefore distinguish a fixed shortcut or automation rule, a draft-oriented assistant, and an adaptive agent; time saved by composing a draft cannot be treated as evidence for broader authorization.

Authority should rise through a productivity effect ladder: observe and search, create a private draft, make a reversible local edit, communicate or change a shared object, and finally perform financial, account-level, public, or irreversible actions. Identity providers and application services--not the language model--should enforce account scope, recipient and domain restrictions, sharing visibility, spending and rate limits, confirmation requirements, idempotency, and audit retention. Untrusted content inside a webpage, message, attachment, or document cannot grant permission, widen scope, or override standing user policy. Ambiguous recipients, inaccessible source objects, expired authentication, conflicting instructions, or a materially changed price or schedule require clarification or handoff.

Evaluation should compare macros, fixed workflows, draft assistants, and adaptive agents under identical tasks, permissions, and budgets. Measure end state and collateral effects, recipient and audience errors, authorization, confirmation, preference misuse, duplication, recovery, corrections, workload, latency, and cost under stale, concurrent, adversarial, and interrupted conditions. The invariant is \emph{intent-preserving effect}: persistent effects require a current workflow case, application scope, exact preview-bound consent when required, and a verified receipt; preferences never substitute for authority.

\subsection{IT operations and cybersecurity}

IT operations makes the paper's central distinction unusually sharp. An agent may summarize an alert, correlate logs, retrieve a runbook, or propose a command, but none of these outputs proves that it has identified the affected system, distinguished cause from symptom, or selected a safe intervention. Operational evidence is heterogeneous, delayed, incomplete, and sometimes adversarial; operational credentials can change production at machine speed. The system must therefore transform telemetry into a bounded remediation case rather than let a plausible narrative flow directly into a shell, cloud, identity, or security control action.

The persistent \emph{operational incident case} binds incident, service, environment, severity, impact, owner, configuration, and policy versions to provenance-bearing telemetry. Observations remain separate from causal hypotheses and falsifying tests; runbooks and threat intelligence remain evidence rather than instructions; credentials are represented only by principal, scope, and expiry.

A proposed remediation specifies target, arguments, preconditions, expected observables, blast radius, privilege, reversibility, idempotency, approval, and rollback before becoming a command. Attempts, diffs, receipts, verification queries, and observation windows append to the case. An exit code proves execution, not recovered service, removed persistence, consistent data, or tenant isolation.

The incident case moves through an explicit state machine: detect, triage, scope, hypothesize, gather and test evidence, propose, simulate or canary, approve, execute, verify, observe, and resolve, rollback, reopen, or escalate. NIST SP 800-61 Rev.~3 similarly places preparation, detection, response, and recovery within continuing cybersecurity risk management rather than treating response as a single reactive step~\cite{nist2025incidentresponse}. State must be refreshed immediately before execution because deployments, traffic, credentials, incident ownership, and concurrent operator actions can change. Resolution requires both an authoritative configuration or security state check and stable service postconditions over a declared observation window; a dashboard that briefly turns green is not closure.

Operational benchmarks cover complementary slices of this contract. AIOpsLab deploys microservice environments, injects faults, generates workloads, exports telemetry, and exposes an agent--cloud interface, enabling separate assessment of detection, localization, diagnosis, and mitigation under controlled cloud failures~\cite{chen2025aiopslab}. ITBench broadens the application scope to 94 scenarios derived for site reliability, compliance and security, and financial operations work. In its reported evaluation, agents resolved 13.8\% of SRE scenarios, 25.2\% of compliance and security scenarios, and none of the FinOps scenarios~\cite{jha2025itbench}. These figures are not directly comparable with AIOpsLab because the tasks, environments, agents, and success rules differ. Their shared lesson is narrower: recognizing a fault or producing a technically relevant step does not establish end-to-end resolution.

Open-ended security investigation exposes an even larger evidence coverage problem. The Cyber Defense Benchmark places agents over 75,000--135,000 Windows event records per episode and asks them to hunt across 26 campaigns covering 105 recorded attack procedures. The strongest of five evaluated models flagged only 3.8\% of malicious events on average, and no run found every malicious event~\cite{chona2026cyberdefense}. SecRespond begins after compromise: 10 ranges span four entry point types, 21 MITRE ATT\&CK techniques, and five operating systems, and agents must analyze forensic disk state, reconstruct intrusions, and prepare remediation. Across 23 model configurations, no model completed detection and remediation on any range; agents more readily surfaced alert-exposed problems than silent persistence or a complete verified response~\cite{wang2026secrespond}. The suites use different artifacts and graders, but both show why an incident case needs explicit coverage obligations and unresolved hypotheses instead of treating the first plausible finding as the incident.

The evidence channel is itself an attack surface. Logs, tickets, webpages, emails, threat reports, package metadata, and even retrieved runbooks can contain attacker-controlled text. Agent Security Bench spans ten scenarios, more than 400 tools, and attacks against prompts, tool use, and memory; its highest reported average attack success rate exceeded 84.3\% in the evaluated configurations~\cite{zhang2024agentsecuritybench}. AgentDojo separately couples 97 benign tasks to 629 security test cases involving untrusted tool data~\cite{debenedetti2024agentdojo}. Neither benchmark estimates the probability of compromise in a particular production operation. They demonstrate that delimiter prompts and model instructions are not authority boundaries. Provenance and trust labels must survive interpretation, and untrusted content must be unable to widen capabilities, select a privileged target, alter an approved change, or suppress required verification.

The current evidence base remains benchmark-heavy. Fault-injected environments, frozen forensic images, and authored adversarial cases reveal failure modes, but they do not establish lower outage duration, fewer false containments, or acceptable operator workload in deployed organizations. A production claim requires field evidence that binds the exact agent configuration and authority level to realized service and security outcomes. Until such evidence is available, evaluation should distinguish read-only triage assistance, recommendation, supervised remediation, and autonomous execution rather than treating them as one capability. The safest useful system may reduce evidence gathering time while leaving the production decision with an accountable operator.

Authority should therefore rise through an operational effect ladder: read telemetry; correlate and recommend; prepare a change and rollback plan; execute in a sandbox or replay; run a narrow canary; perform a reversible, resource-scoped production change; and only then consider destructive, cross-tenant, identity, network, or data plane actions. Infrastructure and policy services--not the language model--should enforce environment and tenant scope, role-based access, separation of duties, change freezes, concurrency controls, rate and resource limits, break-glass procedure, approval, and credential expiry. High-impact actions require an independent reviewer and a tested rollback path. Missing ownership, ambiguous targets, incomplete telemetry, unavailable verification, scope growth, or conflicting operator changes should halt execution and trigger clarification or escalation.

Evaluation should compare fixed automation, an operations copilot, and an adaptive agent under identical incidents, telemetry, privileges, and budgets. Measure detection, diagnosis, containment errors, recovery time, service loss, security restoration, collateral effects, privilege use, rollback, recurrence, workload, latency, and cost under telemetry, concurrency, credential, command, and rollback faults. The invariant is \emph{verified bounded remediation}: production mutation requires versioned evidence for target, preconditions, authority, effect, blast radius, and recovery, and closes only on authoritative postconditions.

\subsection{Scientific discovery and autonomous laboratories}

Scientific agency begins when a system can choose the next analysis or experiment, not when it can write a plausible hypothesis. The distinction matters because a fluent explanation, an improved objective value, a completed protocol, and a scientifically supported claim are different artifacts. Scientific evidence carries units, uncertainty, conditions, sample identity, instrument state, and alternative explanations. A credible agent must preserve those relationships while converting prior evidence into a bounded experimental action; otherwise an attractive narrative can outrun both the measurement and the domain in which it is valid.

The persistent \emph{scientific investigation case} links the question, competing hypotheses, predicted observations, and claim scope to versioned literature, data, simulations, prior experiments, samples, custody, instruments, and conditions. Observations remain distinct from derived quantities and interpretations so a generated summary cannot become experimental fact.

Each experiment is a typed protocol binding variables, controls, units, samples, calibration, hazards, falsifying outcomes, stopping, analysis, instrument state, material, authority, and abort conditions. Execution appends commands and receipts; measurement appends raw data, calibration, uncertainty, and quality control; analysis appends code, artifacts, and claim status changes. A transcript cannot safely carry this state.

The investigation proceeds as a state machine: frame the question, assemble and assess evidence, formulate competing hypotheses, design the protocol, review feasibility and hazards, simulate or dry-run where possible, authorize, execute, acquire measurements, perform quality control, analyze, update the hypotheses, and close, revise, or escalate. Failed, null, censored, and ambiguous outcomes remain first-class states. The system rechecks sample, instrument, material inventory, calibration, and safety state immediately before a physical action. It does not start the next experiment merely because the last command returned successfully, and it does not close a claim until the measurements and analysis satisfy the protocol's declared decision conditions.

FunSearch illustrates the unusually favorable case in which candidate actions and evaluation can be tightly specified. It combines an LLM that proposes programs with a systematic evaluator that executes and scores them, retaining high-performing programs to guide search. The reported system discovered improved constructions for the cap set problem and new online bin packing heuristics~\cite{romeraparedes2024funsearch}. This is strong evidence that generative search can produce verifiable advances when the candidate is executable and the property of interest is objective and comparatively cheap to test. It does not imply that every scientific objective can be reduced to the evaluator: the score certifies the encoded property under the tested conditions, not every interpretation, downstream consequence, or unrepresented constraint.

Controlled scientific agent benchmarks expose different gaps. ScienceAgentBench derives 102 data analysis tasks from 44 peer-reviewed papers in four disciplines, validated by nine subject-matter experts, and evaluates executable Python artifacts. With three attempts per task, its best reported agent solved 32.4\% independently and 34.3\% with expert-provided knowledge~\cite{chen2025scienceagentbench}. DiscoveryWorld instead requires agents to form hypotheses, manipulate a simulated environment, run experiments, and explain discoveries across 24 tasks, three difficulty levels, and parametric variants; it scores task completion, procedural progress, and explanatory knowledge~\cite{jansen2024discoveryworld}. Their scores are not directly comparable: one tests data-driven coding against task-specific outputs, whereas the other tests long-horizon discovery in a low-fidelity fictional world. Together they show why end-to-end claims should be decomposed into evidence acquisition, experimental action, analysis, and explanation.

Physical closed loops provide stronger evidence about action but remain bounded by their apparatus and measurement contract. The A-Lab combined literature-derived synthesis recipes, active learning, robotics, and X-ray diffraction in 353 experiments over 17 days. In the corrected record, manual reanalysis confirmed 36 of 57 target materials while four earlier reported successes remained inconclusive from diffraction alone~\cite{szymanski2023alab,szymanski2026alabcorrection}. The correction also clarified that ``new'' referred to targets new to the prediction platform rather than necessarily new to science. This history does not erase the demonstrated autonomous synthesis loop. It makes the survey's point concrete: automated classification can guide the next experiment, while the final scientific claim must retain measurement ambiguity, human reanalysis, and its precise novelty scope.

Coscientist connects web and documentation search, code execution, and laboratory APIs, and was demonstrated across six task classes including instrument control and optimization of palladium-catalysed cross-couplings~\cite{boiko2023coscientist}. ChemCrow integrates 18 expert-designed chemistry tools, planned and executed syntheses of an insect repellent and three organocatalysts, and guided a chromophore experiment~\cite{bran2024chemcrow}. In ChemCrow's comparison, expert chemists favored the tool-augmented system on chemical correctness and task completion, whereas a GPT-4 evaluator favored fluent but factually weaker baseline answers. These demonstrations establish that specialized tools and instrumentation can ground parts of chemical work; they also show why another unconstrained language model is not an adequate substitute for domain measurements and expert judgment.

More open-ended systems, including AI Scientist and SciAgents, explore automated experiment cycles, manuscript production, hypothesis generation, and multi-agent reasoning~\cite{lu2024aiscientist,ghafarollahi2024sciagents}. Their artifacts can be useful proposals, but proposal volume, simulated novelty, or language model review does not establish a scientific result. The evidence base therefore supports a graded claim: agents can search, analyze, optimize, and sometimes operate a tightly bounded experimental loop, while general autonomous discovery across changing instruments, hazards, evidence standards, and scientific communities remains unproven.

Authority should rise through a scientific effect ladder: retrieve and organize literature; analyze an existing dataset; propose a hypothesis; prepare code or a protocol; run a simulation; perform a dry-run or low-risk instrument check; execute a bounded experiment; and only then approach actions involving hazardous chemicals, biological agents, high energy, expensive or scarce samples, or irreversible transformations. Laboratory scheduling, inventory, instrument controllers, and safety systems--not the language model--should enforce material identity, facility access, parameter envelopes, interlocks, waste and exposure limits, resource budgets, and emergency stops. Independent review is required when the hazard, uncertainty, cost, or scientific consequence exceeds the case's approved scope.

Evaluation should compare assistants, fixed experimental pipelines, and adaptive agents under identical questions, evidence, apparatus, materials, and budgets. Measure hypothesis discrimination, protocol and control validity, unit and sample integrity, calibration, safety, measurement quality, uncertainty, claim support, expert correction, time, and cost under conflicting, mislabeled, drifting, missing, and ambiguous conditions. The invariant is \emph{measurement-bounded discovery}: action requires a complete protocol, apparatus state, hazards, authority, controls, and abort path; claims cannot exceed identified measurements and quality checks.

\subsection{High-stakes professional support}

High-stakes professional support is not a single application category; it is a common transition from evidence to an institutionally consequential decision. The same generated paragraph can be a harmless draft, a clinician's working hypothesis, a legal position, an investment recommendation, or the stated reason for denying a public benefit. A system becomes materially agentic when it selects evidence, advances a case, or invokes a tool that changes another person's care, rights, money, or access to services. The governing question is therefore not whether the output sounds expert, but whether every case transition remains evidentially supported, within delegated authority, reviewable, and contestable.

The shared \emph{professional case record} separates assertions, verified and contested facts, model suggestions, professional judgments, authorized decisions, and effects while preserving provenance, applicability, confidentiality, consent, deadlines, uncertainty, reasons, and review rights. Clinical, legal, financial, and public service cases extend this grammar with domain identities, rules, evidence, constraints, and accountable professionals. Each proposed commitment binds evidence, applicable standard, uncertainty, harms, decision rights, reviewer, executable action, and correction or appeal. Missing identity, consent, authority, current rules, material evidence, or a qualified reviewer is a stop state.

Healthcare evidence illustrates both promise and claim limits. AMIE was compared with 20 primary care physicians in a randomized, double-blind crossover study using 159 simulated, text-based diagnostic consultations. Specialist physicians rated AMIE superior on 30 of 32 evaluation axes and patient actors on 25 of 26 axes~\cite{tu2025amie}. The study required active information gathering rather than diagnosis from a fixed vignette, but its authors explicitly caution that the synchronous text setting was not representative of usual clinical practice. MDAgents dynamically selected solo or group collaboration and achieved the best result on seven of ten medical benchmarks, with gains of up to 4.2\% over prior best methods~\cite{kim2024mdagents}. These studies support narrower claims about simulated dialogue, medical knowledge, and benchmark reasoning. They do not establish patient outcomes, workflow safety, or authority for an agent to diagnose, prescribe, order, or communicate a final care decision without a responsible clinician and a current clinical record.

Legal evidence exposes a similar progression. LegalBench contains 162 expert-constructed tasks across six forms of legal reasoning and evaluated 20 models~\cite{guha2023legalbench}. LegalAgentBench adds 300 multi-hop and writing tasks, 17 corpora, and 37 retrieval, database, and calculation tools in the Chinese legal domain, while recording progress through intermediate tool use~\cite{li2024legalagentbench}. The latter is closer to agentic research, but a correct benchmark answer is not a valid position in a live matter with changing law, disputed facts, privilege, deadlines, and client duties. In a preregistered evaluation of commercial retrieval-augmented legal research products, Magesh et al. found hallucinations in 17--33\% of responses, despite rates below those of a general-purpose model~\cite{magesh2025legalrag}. A linked source is therefore not enough: the case must bind each proposition to the authority's holding, court, jurisdiction, date, subsequent treatment, and applicability to the client's facts.

Financial benchmarks likewise distinguish analytic capability from delegated fiduciary or transactional authority. FinBen spans 42 datasets and 24 tasks across eight areas, including information extraction, question answering, risk management, forecasting, and decision-making. Its evaluation of 21 models found stronger performance on extraction and textual analysis than on advanced reasoning, generation, and forecasting~\cite{xie2024finben}. FinCon organizes analyst agents under a manager, updates investment beliefs through risk-controlled reflection, and evaluates historical stock trading and portfolio management trajectories against LLM, reinforcement learning, buy-and-hold, and classical portfolio baselines using return, Sharpe ratio, and drawdown~\cite{yu2024fincon}. Such comparisons can establish performance under a declared dataset and backtest protocol. They do not establish suitability, robustness to regime change, transaction execution, market impact, compliance, or realized client benefit. Those claims require a case record that freezes data vintage, mandate, holdings, exposure, constraints, approvals, orders, fills, and reconciliation.

Public administration supplies stronger field evidence about the consequences of weak case structure, although not from an LLM agent. The Royal Commission into Australia's Robodebt Scheme recommended documenting the sources and bases of assumptions, obtaining legal review of end-to-end data exchanges, retaining vulnerability information on recipient files, providing a clear path to review, disclosing automated decision-making, and making business rules available for independent scrutiny~\cite{robodebt2023report}. This is important counterevidence: automation does not become legitimate because a person can inspect its final message. The evidentiary basis, applicable law, individual circumstances, reasons, notice, and route to challenge must travel with each decision, and recovery activity must pause when the underlying case is disputed or under review.

Authority should rise through a professional effect ladder: retrieve and organize sources; identify missing evidence; prepare a chronology, calculation, or draft; generate options and counterarguments; recommend an action; obtain accountable professional approval; and only then perform a narrow, structured, reversible administrative step. Diagnosis and treatment, legal filing or representation, investment and credit decisions, benefit eligibility, debt recovery, and enforcement should not be delegated to an unconstrained language model. Clinical, matter management, trading, case management, and policy systems--not the model--should enforce identity and consent, licensure and privilege, jurisdiction, suitability, position and loss limits, separation of duties, deadlines, notice and appeal, approval scope, and audit retention.

Evaluation should compare existing professional workflows, drafting copilots, and adaptive agents under identical cases, evidence, authority, time, and budgets. Measure evidence and attribution, currency, uncertainty, missed conflicts, professional correction, workload, disparate effects, inappropriate decisions, privacy, notice and appeal, outcomes, recovery, latency, and cost under identity, authority, policy, consent, and adversarial failures. The invariant is \emph{professionally accountable case progression}: consequential action requires a current case, applicable rules, decision rights, a responsible professional, review path, and expected effect; fluency never creates professional authority.

\subsection{Cross-domain deployment patterns}

Across domains, the nouns change but the control problem does not. A research assistant works with claims and sources; a support agent with customers, policies, and transactions; a coding agent with requirements, diffs, and tests; an operations agent with incidents and infrastructure; a laboratory agent with hypotheses, samples, and measurements; and a professional system with cases, duties, and decisions. In each setting, reliability depends on preserving the relation between what was observed, what was inferred, what was proposed, what was authorized, what was executed, and what happened afterward. Long-horizon agency is therefore not simply a longer conversation or a larger number of tool calls. It is a sequence of durable, inspectable state transitions.

The shared grammar is small: evidence items preserve provenance and applicability; case state preserves facts, disputes, obligations, and authority; proposals bind effects to evidence and validation; commitments freeze target, scope, approval, and recovery; receipts record authoritative postconditions and collateral effects. Summaries may be regenerated, but admitted evidence, approvals, and receipts cannot be silently rewritten.

The shared lifecycle is evidence acquisition, case update, bounded proposal, validation, authority, controlled commitment, and reconciliation. Models interpret ambiguity and select among permitted steps; deterministic services and accountable people control admission, authority, and effects. Missing evidence or authority produces an unresolved or escalated state. Six deductions follow:
\begin{enumerate}
    \item \textbf{Begin with the narrowest architecture.} Fixed retrieval and workflows should handle specifiable transitions. Agentless remained competitive on SWE-bench Lite, while field support evidence showed value from suggestions without transaction authority~\cite{xia2024agentless,brynjolfsson2025generativeaiatwork}. Compare fixed and adaptive designs under the same model, tools, and budget.
    \item \textbf{Center the authoritative artifact.} Cases close against database state, accepted diffs, application state, incident postconditions, measurements, or accountable decisions. AppWorld and the A-Lab correction show why narrative cannot substitute for state or measurement~\cite{trivedi2024appworld,szymanski2026alabcorrection}.
    \item \textbf{Separate proposal from commitment.} Proposals are editable; commitments freeze an exact action, target, authority, validity window, and recovery path. ClawsBench and workplace studies show that completion, usefulness, and permission are distinct~\cite{li2026clawsbench,brynjolfsson2025generativeaiatwork,dillon2025shiftingwork}.
    \item \textbf{Use domain-specific validators.} FunSearch benefits from executable scoring, whereas ChemCrow and legal RAG evidence expose limits of fluent judges and linked citations~\cite{romeraparedes2024funsearch,bran2024chemcrow,magesh2025legalrag}. Validators judge frozen proposals without rewriting criteria or enlarging authority.
    \item \textbf{Preserve provenance.} ALCE exposes incomplete citation support, and Agent Security Bench shows that evidence channels can carry attacks~\cite{gao2023alce,zhang2024agentsecuritybench}. Evidence may inform interpretation but cannot grant capability, change approval, or suppress verification.
    \item \textbf{Evaluate the human workflow.} Compare manual work, fixed automation, assistants, proposal-only agents, and action-capable agents under common authority. Field productivity evidence and AMIE's simulated results support different claims~\cite{brynjolfsson2025generativeaiatwork,cui2026highskilledwork,tu2025amie}; measure accepted outcomes, correction, workload, harms, and accountability.
\end{enumerate}

These patterns form three planes: evidence stores the case, decision interprets ambiguity and proposes permitted steps, and effect enforces identity, policy, approval, recovery, and authoritative receipts. Autonomy should therefore be reported as a contract over admitted evidence, selectable transitions, maximum effect, validation, recovery, and accountable authority rather than as a scalar.

The cross-domain invariant is \emph{representation-bounded action}: proposals require typed evidence and criteria; commitments freeze target, preconditions, authority, limits, and recovery; cases close only on authoritative postconditions. Missing or contradictory elements trigger retrieval, revision, withholding, rollback, reopening, or escalation. The evidence supports a stronger conclusion than ``more autonomy is better'': agentic value comes from turning heterogeneous evidence into bounded, inspectable effects.
